\subsection{Matériel expérimental}
L'étude expérimentale du microscope repose sur l'utilisation des éléments suivants :

\begin{itemize}
    \item \textbf{Microscope optique} classique. 
    \item \textbf{Éclairage} additionnel assuré par une lampe électrique.
    \item \textbf{Oculaires} :
    \begin{itemize}
        \item Un oculaire classique ;
        \item Un oculaire micrométrique.
    \end{itemize}
    \item \textbf{Dispositifs de mesure et de visée} :
    \begin{itemize}
        \item Mire de Foucault ;
        \item Lame micrométrique (micromètre-objet) graduée au $1/100$ mm ;
        \item Chambre claire (système de projection à miroirs) ;
        \item Règle graduée.
    \end{itemize}
    \item \textbf{Supports et objets d'étude} :
    \begin{itemize}
        \item Diapositive contenant du papier millimétré (pour l'étude du cercle oculaire) ;
        \item Préparation biologique (lame de larves de Pluteus).
    \end{itemize}
\end{itemize} 

\subsection{Protocoles expérimentaux}

\subsubsection{Détermination de la limite de séparation de l'œil}

Cette manipulation vise à mesurer l'angle minimal sous lequel l'œil humain peut distinguer deux points distincts. Le protocole suivi est le suivant :

\begin{itemize}
    \item Éclairer de manière intense la mire de Foucault.
    \item Reculer progressivement par rapport à la mire jusqu'à ce que les traits noirs et blancs ne soient plus discernables, l'image fusionnant en un gris uniforme.
    \item Mesurer la distance limite $D_m$ séparant l'œil de la mire à cet instant précis.
    \item Relever la période $a$ de la mire (distance entre deux bandes noires ou deux bandes blanches successives) pour calculer l'angle limite de séparation $\epsilon = \frac{a}{D_m}$.
\end{itemize}

\subsubsection{Mise au point et observation}

Cette étape permet d'obtenir une image nette de l'objet et de déterminer les conditions de visée pour chaque configuration optique.

\begin{itemize}
    \item Sélectionner l'objectif de plus faible grandissement ($\times 4$).
    \item Positionner le tube optique à sa distance minimale de la préparation par observation latérale.
    \item Remonter progressivement le tube à l'aide du bouton de crémaillère jusqu'à l'apparition de l'image dans l'oculaire.
    \item Ajuster la netteté finale via la vis micrométrique.
    \item Mesurer et noter :
    \begin{itemize}
        \item La distance entre la sortie de l'oculaire et l'objet ;
        \item La distance estimée entre l'œil de l'observateur et l'objet.
    \end{itemize}
    \item Répéter l'ensemble de ces opérations et mesures pour les deux autres objectifs ($\times 10$ et $\times 40$).
\end{itemize}

\subsubsection{Caractérisation du cercle oculaire}

Cette manipulation permet de localiser et de mesurer la section minimale du faisceau lumineux sortant de l'instrument, correspondant à l'image de la pupille d'entrée par l'oculaire.

\begin{itemize}
    \item Réaliser la mise au point sur le micromètre-objet en utilisant l'objectif $\times 10$.
    \item Placer la diapositive de papier millimétré à la sortie de l'oculaire.
    \item Déplacer verticalement la diapositive le long de l'axe optique pour identifier la position de section minimale du disque lumineux.
    \item Mesurer le diamètre de ce cercle oculaire à l'aide des graduations du papier millimétré.
    \item Relever la distance séparant le cercle oculaire du verre de la lentille de sortie de l'oculaire.
\end{itemize}

\subsubsection{Mesure du grossissement commercial et calcul de la puissance intrinsèque}

Cette manipulation permet de déterminer expérimentalement le grossissement $G_c$ du microscope en comparant la taille de l'image perçue à celle de l'objet réel.

\begin{itemize}
    \item Installer la chambre claire sur l'oculaire du microscope.
    \item Placer une règle graduée verticalement à une distance de $25$ cm de l'œil, en veillant à ce qu'elle soit suffisamment éclairée pour être visible par réflexion.
    \item Observer simultanément le micromètre-objet (par transmission) et la règle graduée (par réflexion via le miroir de la chambre claire).
    \item Faire coïncider les deux images et relever la longueur $L$ mesurée sur la règle correspondant à un nombre $n$ de divisions du micromètre-objet.
    \item Calculer le grossissement commercial et la puissance intrinsèque en utilisant la relation entre la taille de l'image observée 
    sur la règle et la taille réelle de l'objet ($1 \mathrm{ div} = 0,01$ mm).
\end{itemize}

\subsubsection{Mesure du champ latéral}

Cette manipulation permet de déterminer l'étendue circulaire de l'espace visible à travers l'instrument pour une configuration donnée.

\begin{itemize}
    \item Réaliser la mise au point sur le micromètre-objet en utilisant l'objectif $\times 40$.
    \item Aligner le zéro de la graduation du micromètre-objet avec l'un des bords du champ visuel circulaire.
    \item Dénombrer le nombre total de divisions visibles sur toute l'étendue du diamètre du champ.
    \item Convertir ce nombre de graduations en distance réelle afin de déterminer le diamètre du champ latéral.
\end{itemize}

\subsubsection{Étalonnage et mesure d'objet (Pluteus)}

Cette manipulation finale consiste à étalonner l'instrument pour effectuer des mesures biométriques précises sur un échantillon biologique.

\paragraph{Procédure d'étalonnage}
\begin{itemize}
    \item Substituer l'oculaire classique par l'oculaire micrométrique et noter son grossissement.
    \item Placer le micromètre-objet sur la platine et réaliser la mise au point avec l'objectif $\times 10$.
    \item Faire pivoter le tube porte-oculaire afin de rendre les graduations de l'oculaire et celles du micromètre-objet parfaitement parallèles.
    \item Établir la correspondance entre les divisions de l'oculaire et les dimensions connues du micromètre-objet ($1/100$ mm).
    \item Déterminer précisément la valeur réelle d'une graduation de l'oculaire pour cette configuration.
\end{itemize}

\paragraph{Procédure de mesure} 
\begin{itemize}
    \item Remplacer le micromètre-objet par la lame de Pluteus.
    \item Sélectionner une larve complète et de taille représentative, puis l'amener au centre du champ visuel.
    \item Effectuer la mise au point avec l'objectif $\times 10$.
    \item Mesurer les dimensions caractéristiques de la larve (longueur et largeur) en dénombrant les graduations oculaires correspondantes.
    \item Convertir ces mesures en dimensions réelles à l'aide du facteur d'étalonnage calculé précédemment.
\end{itemize}


